\section{Conclusion}\label{sec:conclusion}

We distilled small students from teachers that mix a reasoning model and an instruct model with ratio $\lambda$, and examined what changes in the student's LoRA adapter when the teacher stops thinking. The student's outputs followed the teacher, with reflective markers and length shrinking together, so they could not show what had changed. The adapter, by contrast, sharply changes the directions along which it updates the weights at the $\lambda$ where its own teacher stops thinking, in all three pairs, and the adapters before that point always share the same small region.

This shared region acts as a switch that turns thinking on and off. Removing it turns thinking off and lets the student finish short solutions, unlike a length-matched cut, whereas removing the same energy from other directions outside the shared region does not turn thinking off to the same degree. Amplifying it with the norm held fixed does not increase thinking, and grafting it into an adapter that gave short answers turns thinking on but leaves the student unable to finish its solutions. Reading the student's internal thinking features leads to mostly the same conclusions. The shared region also appears beyond the training problems and in distillation without interpolation: removing the region found only from MATH turns thinking off on GSM8K and GPQA-Diamond, and a little over half of it is captured by just two reasoning-distillation adapters trained without interpolated teachers.

These findings suggest that reasoning distillation into small students stores the teacher's way of thinking separately in a small region of the weights, and that this region only turns thinking on without carrying the solution through to the end; what is needed to finish a solution is not transferred by this region alone. More broadly, a change that cannot be told apart in the student's outputs, where thinking and length collapse together, appears separately in the weights gathered by LoRA, aligned with the $\lambda$ at which the teacher stops thinking.
